\begin{abstract}
\noindent
Se desarrolló una red neuronal de clasificación binaria para predecir la alta letalidad (dos o más fallecidos) de los siniestros viales fatales del Perú, usando la base oficial del Observatorio Nacional de Seguridad Vial (ONSV) 2021--2025 (9\,104 siniestros fatales útiles, 10.2\,\% de casos multifatales). Se ejecutó un pipeline completo: análisis exploratorio, preprocesamiento anti-fuga, ingeniería de 116 características pre-impacto (zona, red vial, clima, curvatura, perfil, superficie, coordenadas y derivadas temporales), entrenamiento de un perceptrón multicapa comparado contra tres líneas base, calibración del umbral de decisión en validación y calibración probabilística post-hoc. El modelo final alcanzó ROC-AUC de 0.7474 y recall-multifatal de 0.8849 en el conjunto de prueba, con intervalos de confianza bootstrap al 95\,\%. El análisis SHAP identificó la clase de siniestro (atropello, despiste, choque) y la zona rural/urbana como los factores de mayor influencia, coherentes con el análisis exploratorio (tasa multifatal rural 15.1\,\% vs.\ urbana 3.6\,\%). La calibración isotónica redujo el Brier de 0.2139 a 0.0872, por lo que la interfaz gráfica en Streamlit muestra probabilidades calibradas reales. Se reportan con honestidad las limitaciones: la regresión logística compite de cerca con la red en F1, el registro 2025 es preliminar y el modelo es una estimación académica de factores de riesgo, no un predictor operacional.
\end{abstract}

\bigskip